\section{Simulation}\label{sec:simulation}

\subsection{Design}

Data were simulated from a joint model with a random intercept and slope,
longitudinal mean $\beta_0 + \beta_1 t + \beta_{\mathrm{age}}\mathrm{age}_i
+ \beta_{\mathrm{sex}}\mathrm{sex}_i$ with $\beta_0 = 2$ and $\beta_1 =
0.5$, $\mathrm{age}_i \sim \N(0, 10^2)$ and $\mathrm{sex}_i \sim
\mathrm{Bernoulli}(0.5)$, $\sigma_b = (0.8, 0.3)$, $\sigma_e = 0.3$, and a
Weibull hazard (shape 1.3) with $\alpha = 0.5$ and a treatment effect
$\gamma = -0.3$, calibrated to an event rate of about one half. Visits
were scheduled at regular intervals up to $t = 10$ and stopped at the event
or censoring time. All fits used the spline baseline hazard of
Section~\ref{sec:setup}, so the fitted model does not contain the true
hazard exactly. Each design was simulated with three seeds, and every seed
was fitted twice on the same data: unrotated, and with the rotation and
its dense block. Both fits report the same parameters, $\beta$ included,
so the comparison is on the quantities a user reports. The design is
paired: both arms see the same data. Effective sample size (ESS) is
NumPyro's (version 0.15) multi-chain autocorrelation estimator, truncated
by Geyer's initial monotone sequence and without rank normalization
\citep{geyer1992,gelman2013}, and $\hat R$ is its split-chain version. ESS
per second divides by the wall time of the whole MCMC run, warm-up and
compilation included; ESS per gradient divides by the number of leapfrog
steps in the sampling phase. Ratios are rotated over unrotated, paired by
seed, with the geometric mean over seeds reported for each design and the
per-seed range in the text.

The core grid crossed about 9 against about 17 visits per subject with
$\rho \in \{0, 0.3, 0.6, 0.9\}$ at $N = 300$, and added $N \in \{150,
600\}$ at 9 visits and $\rho = 0.3$: ten designs, 2 chains of 500 warm-up
and 1000 sampling iterations.

\subsection{Efficiency}

\begin{table}[t]
\centering
\caption{Core grid: ESS per second for the intercept and the time slope
(means over 3 seeds), and the rotated/unrotated ratio (geometric mean of
the paired per-seed ratios).}
\label{tab:core}
\small
\begin{tabular}{rrr rrr rrr r}
\toprule
& & & \multicolumn{3}{c}{intercept} & \multicolumn{3}{c}{time slope} & $\alpha$ \\
\cmidrule(lr){4-6}\cmidrule(lr){7-9}\cmidrule(lr){10-10}
visits & $\rho$ & $N$ & unrot. & rotated & ratio & unrot. & rotated & ratio & ratio \\
\midrule
9  & 0.0 & 300 & 25.9 & 354.4 & 14.2 &  9.1 & 304.9 & 36.4 & 0.85 \\
17 & 0.0 & 300 & 12.1 & 229.4 & 19.3 &  4.0 & 234.9 & 59.2 & 1.32 \\
9  & 0.3 & 300 & 30.0 & 253.3 &  8.6 & 19.0 & 251.1 & 13.4 & 0.72 \\
17 & 0.3 & 300 & 11.3 & 248.5 & 22.2 &  6.4 & 262.5 & 41.9 & 1.33 \\
9  & 0.6 & 300 & 26.1 & 245.7 &  9.4 & 29.9 & 229.3 &  7.7 & 1.00 \\
17 & 0.6 & 300 &  9.6 & 231.6 & 24.0 & 10.0 & 251.2 & 24.9 & 1.59 \\
9  & 0.9 & 300 & 20.4 & 220.7 & 11.1 & 20.3 & 250.3 & 12.8 & 0.90 \\
17 & 0.9 & 300 &  8.9 & 218.9 & 24.9 &  8.6 & 217.7 & 25.6 & 1.59 \\
9  & 0.3 & 150 & 41.8 & 323.7 &  7.6 & 32.6 & 420.4 & 13.1 & 0.76 \\
9  & 0.3 & 600 & 15.9 & 132.0 &  8.2 &  8.8 & 151.7 & 17.7 & 0.77 \\
\bottomrule
\end{tabular}
\end{table}

Table~\ref{tab:core} gives the results. Over all 30 paired fits the
geometric-mean gain was 13.5 for the intercept (per seed 4.5--31.6) and
21.1 for the time slope (5.1--70.8). Wall time was essentially unchanged
(ratio 1.03, range 0.83--1.10), and the rotated fit took slightly fewer
leapfrog steps per iteration (0.89 times; 0.49--1.02), so per gradient
evaluation the gains were 15.6 for the intercept and 24.3 for the time
slope. Almost all of the gain is in ESS per draw (13.9 and 21.7 times):
under the rotation the median ESS over 2000 draws was about 3,450 for
both coefficients, more than the number of draws, against 180--235
unrotated, that is, about 1.7 against 0.1 ESS per draw. There were no divergences, and the worst $\hat R$ on a reported
coefficient was 1.002 rotated and 1.044 unrotated. Posterior means agreed to a median of 0.027 posterior standard deviations, and the largest paired difference was 2.7 Monte Carlo standard errors over 90 pairs. Posterior standard deviations agreed too: rotated over unrotated, 1.01 for
the intercept (per seed 0.94--1.10) and 0.98 for the time slope
(0.91--1.06), within the Monte Carlo error of the unrotated fit's short
effective samples.

The pattern follows Proposition~\ref{prop:ridge}. The gain is largest with
dense visits, where each subject carries the most information about its
own random effects (19--25 times for the intercept, 25--59 for the slope),
and does not grow with $N$ (7.6 at $N = 150$, 8.2 at $N = 600$), because
the ridge is set by per-subject information, not by the number of
subjects. Figure~\ref{fig:gain} shows all three grids against that
information, $\bar n\,\sigma_0^2/\sigma_e^2$.

The association parameter, which the degeneracy does not involve, is not
untouched: its ratio split by design, 0.72--1.00 in the six designs with
about 9 visits and 1.32--1.59 in the four with about 17 (geometric mean
1.04). We do not have an established explanation. On real data
(Section~\ref{sec:realdata}) the effect on $\alpha$ tracked whether the
degeneracy had been setting NUTS's trajectory length, which is a plausible
reading here too but was not measured.

\begin{figure}[t]
\centering
\includegraphics[width=\textwidth]{fig2_gain.pdf}
\caption{ESS per second of the reported coefficient, rotated over unrotated
fit on the same simulated data, against the per-subject information
$\bar n\,\sigma_0^2/\sigma_e^2$ ($\sigma_0 = 0.8$). Large markers:
geometric mean over three seeds in one design; small: single seeds. Core
and stress grids have a random intercept and slope ($q = 2$); the $q = 1$
grid has a random intercept only. In (b), open triangles are the $q = 1$
time slope, which has no random effect and so nothing to rotate. The $q =
1$ intercept ratios are at the standard run length, where most unrotated
fits had not converged; Section~\ref{sec:q1} gives a comparison between
converged fits.}
\label{fig:gain}
\end{figure}

\subsection{Random intercept only}\label{sec:q1}

The same simulator with $\sigma_{b1} = 0$, fitted with a random intercept
only, gave seven designs: about 9, 17 and 35 visits, $N \in \{100, 300\}$,
and one noisier design with $\sigma_e = 0.8$. The intercept gained 6.5--282
times and the subject-constant coefficient of age 6.6--207 times; the time
slope, which has no absorbable direction when there is no random slope,
moved only through the step size (0.79--3.4 times). At this run length,
however, 17 of the 21 unrotated fits had some parameter with $\hat R$ above
1.05, the intercept's reaching 1.27, while every rotated fit had $\hat R
\leq 1.004$, so these ratios compare a converged fit with unconverged ones.
Posterior means agreed to a median of 0.025 posterior standard deviations
(maximum 0.175).

To obtain a comparison between converged fits, we reran three designs
($N = 300$ with about 9, 17 and 35 visits) and extended the unrotated fit
to 2, 4, 8, 16 and 32 times the standard length, with two seeds taken to
the longest runs. A convergence criterion was fixed in advance: split $\hat
R \leq 1.01$ on every monitored parameter and ESS $\geq 400$ for the
intercept, age, the time slope and $\alpha$. Every rotated fit met it at
the standard length. The unrotated fit met it first at 16 times the
standard length with 9 visits, at 16 and 32 times with 17 visits, and not
by 32 times with 35 visits, where its intercept had an ESS of 304--377 from
64,000 draws with $\hat R = 1.002$. Its ESS per draw for the intercept did
not change with run length: 0.009--0.023, 0.007--0.020 and 0.002--0.008 in
the three designs over runs from 1 to 32 times the standard length, against
2.1--2.5, 0.9--1.9 and 1.1--2.4 for the rotated fit. The per-draw advantage,
of order 100 times, is therefore a property of the stationary chains, not
of short runs. In wall time, the unrotated fit needed 3.5 and 6.5--6.6
times as long as the rotated fit to meet the criterion in the first two
designs, and more than 10--14 times in the third. These ratios are much
smaller than the per-draw ones because the rotated fit's time at the
standard length is mostly warm-up and compilation, and it exceeded the ESS
target by a factor of 5 to 12. Between the first fits of each arm to meet
the criterion, posterior means agreed to a median of 0.02 posterior
standard deviations (maximum 0.11) and standard deviations to within 6\%.

\subsection{Stress grid}\label{sec:stress}

Section~\ref{sec:residual} names two regimes where the rotation should be
weakest: few subjects each carrying a lot of information, and strongly
correlated random effects. A stress grid targeted both, with about 35
visits crossed with $N \in \{50, 100\}$ and $\rho \in \{0.3, 0.95\}$, plus
$\rho = 0.95$ at 9 visits ($N = 100$) and at 17 visits ($N = 300$). The
criterion fixed before the run was that the rotation must not fall below
the unrotated fit on the intercept or time slope in any design. It did not
(Table~\ref{tab:stress}); the lowest design was 5.5 times and the lowest
single seed 3.2 times. Per gradient evaluation the gains were 16.7 for the
intercept and 22.0 for the time slope (geometric means over all designs),
and the rotated fit took 0.81 times as many leapfrog steps per iteration.

\begin{table}[t]
\centering
\caption{Stress grid: ESS-per-second ratios, rotated over unrotated
(geometric means over 3 seeds).}
\label{tab:stress}
\small
\begin{tabular}{rrr rrr}
\toprule
visits & $N$ & $\rho$ & intercept & time slope & $\alpha$ \\
\midrule
35 &  50 & 0.30 &  5.5 & 10.3 & 1.60 \\
35 &  50 & 0.95 & 21.1 & 21.6 & 2.42 \\
35 & 100 & 0.30 & 14.5 & 27.1 & 1.58 \\
35 & 100 & 0.95 & 26.0 & 28.6 & 1.73 \\
 9 & 100 & 0.95 &  7.0 &  7.5 & 0.74 \\
17 & 300 & 0.95 & 24.6 & 29.7 & 1.09 \\
\bottomrule
\end{tabular}
\end{table}

The design with the fewest subjects and the most information per subject
gained least, as the fixed-metric residual predicts. Strong correlation
did not reduce the rotation's advantage: at $\rho = 0.95$ the gains were as
large as or larger than at $\rho = 0.3$, partly because the unrotated fit
degrades faster. Nor did it hurt the rotated fit in absolute terms: with
35 visits its ESS per draw for the intercept was 1.25 at $\rho = 0.95$
against 0.68 at $\rho = 0.3$ with $N = 50$, and 1.50 against 1.18 with
$N = 100$. The predicted
weakness is visible only against a variant that removes the absorbable
directions from the likelihood's coordinates and restores the reported
coefficients with an exact correction (Section~\ref{sec:related}). That
variant has no fixed-metric residual, and in the 35-visit designs it was
1.4--2 times more efficient than the rotation, while being about equal to
it with 9 visits. Posterior means agreed with the unrotated fit's to a
median of 0.029 posterior standard deviations (maximum 0.166, 90 pairs),
there were no divergences, and wall time was 0.86--1.08 times the unrotated fit's (geometric means per design). $\alpha$, which the degeneracy does not involve, fell to 0.74 in the 9-visit design (Table~\ref{tab:stress}).

\subsection{Calibration}

Proposition~\ref{prop:exact} says the rotation cannot change the posterior,
so a calibration study checks the implementation rather than the method.
With $N = 300$ and 100 replicates, every one of which passed an $\hat R
\leq 1.05$ check on the regression parameters (so none was excluded),
coverage of 95\% intervals agreed within
0.01 and the ratio of mean posterior standard deviation to the standard
deviation of posterior means within 0.015 for every estimand
(Table~\ref{tab:cal}); the Monte Carlo standard error of a 95\% coverage
is about 0.011. Two features appear in both arms and belong to the model
and design: slightly wide intervals for $\beta_0$, and an upward bias of
about 3 Monte Carlo standard errors (2.8\%) in $\alpha$. The latter is
consistent with ordinary finite-sample bias. In a separate study
(\texttt{dev/study\_alpha\_bias.R}, Appendix~\ref{app:repro}), neither the
spline baseline hazard nor the simulator's hazard calibration explained
it (paired changes of $-0.0011$
and $-0.0016$, standard error about 0.0025), and at $N = 1200$ it fell to
1.1\%, 1.7 Monte Carlo standard errors, with coverage 0.92.

\begin{table}[t]
\centering
\caption{Calibration, $N = 300$, 100 replicates: coverage of 95\%
intervals, ratio of mean posterior SD to the SD of posterior means, and
bias in Monte Carlo standard errors.}
\label{tab:cal}
\small
\begin{tabular}{l rrr rrr}
\toprule
& \multicolumn{3}{c}{unrotated} & \multicolumn{3}{c}{rotated} \\
\cmidrule(lr){2-4}\cmidrule(lr){5-7}
estimand & cover & SD ratio & bias/se & cover & SD ratio & bias/se \\
\midrule
$\beta_0$           & 0.99 & 1.164 &  1.68 & 0.99 & 1.150 &  1.55 \\
$\beta_{\mathrm{time}}$ & 0.97 & 1.027 & $-0.30$ & 0.96 & 1.038 & $-0.32$ \\
$\alpha$            & 0.93 & 0.922 &  2.97 & 0.93 & 0.917 &  2.97 \\
$\sigma_{b0}$       & 0.94 & 0.914 &  0.63 & 0.95 & 0.906 &  0.68 \\
\bottomrule
\end{tabular}
\end{table}
